\chapter{Tiền bạc và giá trị}
\markboth{Tiền bạc và giá trị}{Money and Value}

\section{Người trẻ kiếm được tiền đầu tiên rồi tiêu để chứng minh mình lớn.}
\begin{truyen}{Người trẻ kiếm được tiền đầu tiên rồi tiêu để chứng minh mình lớn.}{Spending is not maturity.}
\chuhoa{C}{ậu mua những thứ nổi bật nhất để cảm thấy mình đã thành công. Cuối tháng, tài khoản gần như trống rỗng và lòng vẫn không yên. Cậu mới hiểu kiếm được tiền là một chuyện, dùng tiền cho đúng là chuyện khác.}
\textit{(He bought flashy things to feel successful. At the end of the month, his account was nearly empty and his heart still restless. He learned that earning money and using it wisely are different skills.)}
\end{truyen}
\begin{baihoc}
Trưởng thành tài chính không nằm ở chỗ tiêu được bao nhiêu, mà ở chỗ biết giữ cho điều gì.
\textit{(Financial maturity is not about how much you can spend, but about what you know how to keep money for.)}
\end{baihoc}
\begin{ghinhoanh}
\textbf{Short English to remember:}

Spending is not maturity.

Money needs direction.
\end{ghinhoanh}
\begin{tuhocnhanh}
1. Em đang dùng tiền để sống tốt hơn hay để được nhìn tốt hơn? \textit{(Are you using money to live better or to look better?)}

2. Khoản nào trong chi tiêu của em phản ánh sự bốc đồng nhiều nhất? \textit{(Which part of your spending reflects impulsiveness the most?)}
\end{tuhocnhanh}
\ngancach

\section{Một người kiếm khá nhưng luôn thiếu vì muốn bằng bạn bè.}
\begin{truyen}{Một người kiếm khá nhưng luôn thiếu vì muốn bằng bạn bè.}{Comparison eats income.}
\chuhoa{A}{nh không nghèo vì thu nhập quá thấp, mà vì mắt luôn nhìn sang đời người khác. Mỗi lần thấy ai mua gì, anh lại muốn theo. Đến lúc mệt mỏi vì chạy theo, anh mới hiểu so sánh là một cái hố không có đáy.}
\textit{(He was not poor because his income was too low, but because his eyes were always on others. Each time someone bought something, he wanted the same. Only after growing exhausted did he see that comparison is a bottomless hole.)}
\end{truyen}
\begin{baihoc}
Tiền không đủ thường bắt đầu từ lòng không biết đủ.
\textit{(Not having enough money often begins with not having enough contentment.)}
\end{baihoc}
\begin{ghinhoanh}
\textbf{Short English to remember:}

Comparison eats income.

Enough begins inside.
\end{ghinhoanh}
\begin{tuhocnhanh}
1. Em có đang chi tiền để đuổi theo hình ảnh của người khác không? \textit{(Are you spending money to chase other people's image?)}

2. Điều gì với em là đủ? \textit{(What does enough mean to you?)}
\end{tuhocnhanh}
\ngancach

\section{Người đàn ông chỉ lo kiếm tiền mà bỏ quên sức khỏe.}
\begin{truyen}{Người đàn ông chỉ lo kiếm tiền mà bỏ quên sức khỏe.}{Money cannot buy back the lost years.}
\chuhoa{C}{ả tuổi trẻ anh đổi lấy các hợp đồng, những chuyến đi liên tục và những đêm thiếu ngủ. Khi có tiền chữa bệnh, anh mới hiểu tiền có thể mua thuốc, nhưng không mua lại quãng đời khỏe mạnh đã mất.}
\textit{(He traded his youth for contracts, constant travel, and sleepless nights. When he finally had money for treatment, he learned that money can buy medicine, but not the healthy years already gone.)}
\end{truyen}
\begin{baihoc}
Kiếm tiền để sống tốt là đúng; sống tệ đi để kiếm tiền thì thường là hiểu muộn nhất.
\textit{(Earning money to live well makes sense; living badly in order to earn money is a lesson often learned too late.)}
\end{baihoc}
\begin{ghinhoanh}
\textbf{Short English to remember:}

Money cannot buy back lost years.

Protect what money cannot replace.
\end{ghinhoanh}
\begin{tuhocnhanh}
1. Em đang hy sinh điều gì cho tiền bạc? \textit{(What are you sacrificing for money?)}

2. Điều đó có thật sự đáng không? \textit{(Is it truly worth it?)}
\end{tuhocnhanh}
\ngancach

\section{Một người nghĩ giàu lên rồi mới cần sống tử tế.}
\begin{truyen}{Một người nghĩ giàu lên rồi mới cần sống tử tế.}{Character should not wait for wealth.}
\chuhoa{A}{nh nghĩ khi khá giả hơn mình sẽ biết cho đi, biết đỡ người khác, biết sống đàng hoàng hơn. Nhưng tiền chỉ làm rõ con người vốn có; nếu lòng đã hẹp khi nghèo, nhiều khi có tiền lòng vẫn hẹp.}
\textit{(He thought that once he became rich, he would become generous and decent. But money only magnifies what is already there; if the heart is narrow while poor, it may remain narrow while rich.)}
\end{truyen}
\begin{baihoc}
Tiền là công cụ làm lớn lên tính cách, không phải phép màu tạo ra tính cách mới.
\textit{(Money is a tool that enlarges character, not magic that creates new character.)}
\end{baihoc}
\begin{ghinhoanh}
\textbf{Short English to remember:}

Character should not wait for wealth.

Money reveals more than it changes.
\end{ghinhoanh}
\begin{tuhocnhanh}
1. Em đang đợi điều kiện tốt hơn mới sống đúng hơn không? \textit{(Are you waiting for better conditions before living rightly?)}

2. Em muốn tiền làm lớn điều gì trong con người mình? \textit{(What do you want money to enlarge in your character?)}
\end{tuhocnhanh}
\ngancach

\section{Một cô gái đổi sở thích, bạn bè, cả thời gian ngủ vì muốn kiếm nhiều hơn.}
\begin{truyen}{Một cô gái đổi sở thích, bạn bè, cả thời gian ngủ vì muốn kiếm nhiều hơn.}{If money costs your whole self, the price is too high.}
\chuhoa{N}{hìn vào thu nhập tăng lên, ai cũng khen cô giỏi. Chỉ cô biết mình đang dần mất tiếng cười, mất năng lượng, mất cả cảm giác thân thuộc với chính mình. Tới lúc đó cô mới hiểu không phải khoản lời nào cũng là lời.}
\textit{(Everyone praised her rising income. Only she knew she was slowly losing her laughter, energy, and sense of self. Then she understood that not every gain is truly a gain.)}
\end{truyen}
\begin{baihoc}
Nếu một mục tiêu tài chính buộc ta đánh đổi gần hết những gì làm mình thành người, cái giá ấy thường quá đắt.
\textit{(If a financial goal forces us to trade away almost everything that makes us human, the cost is often too high.)}
\end{baihoc}
\begin{ghinhoanh}
\textbf{Short English to remember:}

Not every gain is a gain.

Do not get rich by losing yourself.
\end{ghinhoanh}
\begin{tuhocnhanh}
1. Điều gì trong em đang bị tiền bạc đẩy lùi? \textit{(What part of you is being pushed away by the chase for money?)}

2. Giới hạn nào em không muốn đánh đổi thêm nữa? \textit{(What limit do you no longer want to cross?)}
\end{tuhocnhanh}
\ngancach